%%%%%%%%%%%%%%%%%%%%%%%%%%%%% Define Article %%%%%%%%%%%%%%%%%%%%%%%%%%%%%%%%%%
\documentclass{article}
%%%%%%%%%%%%%%%%%%%%%%%%%%%%%%%%%%%%%%%%%%%%%%%%%%%%%%%%%%%%%%%%%%%%%%%%%%%%%%%

%%%%%%%%%%%%%%%%%%%%%%%%%%%%% Using Packages %%%%%%%%%%%%%%%%%%%%%%%%%%%%%%%%%%
\usepackage[utf8]{inputenc}
\usepackage[spanish]{babel}
\usepackage{geometry}
\usepackage{graphicx}
\usepackage{amssymb}
\usepackage{amsmath}
\usepackage{amsthm}
\usepackage{empheq}
\usepackage{mdframed}
\usepackage{booktabs}
\usepackage{color}
\usepackage{psfrag}
\usepackage{pgfplots}
\usepackage{bm}
\usepackage{hyperref}
%%%%%%%%%%%%%%%%%%%%%%%%%%%%%%%%%%%%%%%%%%%%%%%%%%%%%%%%%%%%%%%%%%%%%%%%%%%%%%%

%%%%%%%%%%%%%%%%%%%%%%%%%% Page Setting %%%%%%%%%%%%%%%%%%%%%%%%%%%%%%%%%%%%%%%
\geometry{a4paper, margin=2.5cm}

%%%%%%%%%%%%%%%%%%%%%%%%%% Define some useful colors %%%%%%%%%%%%%%%%%%%%%%%%%%
\definecolor{ocre}{RGB}{243,102,25}
\definecolor{mygray}{RGB}{243,243,244}
\definecolor{deepGreen}{RGB}{26,111,0}
\definecolor{shallowGreen}{RGB}{235,255,255}
\definecolor{deepBlue}{RGB}{61,124,222}
\definecolor{shallowBlue}{RGB}{235,249,255}
%%%%%%%%%%%%%%%%%%%%%%%%%%%%%%%%%%%%%%%%%%%%%%%%%%%%%%%%%%%%%%%%%%%%%%%%%%%%%%%

%%%%%%%%%%%%%%%%%%%%%%%%%% Define an orangebox command %%%%%%%%%%%%%%%%%%%%%%%%
\newcommand\orangebox[1]{\fcolorbox{ocre}{mygray}{\hspace{1em}#1\hspace{1em}}}
%%%%%%%%%%%%%%%%%%%%%%%%%%%%%%%%%%%%%%%%%%%%%%%%%%%%%%%%%%%%%%%%%%%%%%%%%%%%%%%

%%%%%%%%%%%%%%%%%%%%%%%%%%%% English Environments %%%%%%%%%%%%%%%%%%%%%%%%%%%%%
\newtheoremstyle{mytheoremstyle}{3pt}{3pt}{\normalfont}{0cm}{\rmfamily\bfseries}{}{1em}{{\color{black}\thmname{#1}~\thmnumber{#2}}\thmnote{\,--\,#3}}
\newtheoremstyle{myproblemstyle}{3pt}{3pt}{\normalfont}{0cm}{\rmfamily\bfseries}{}{1em}{{\color{black}\thmname{#1}~\thmnumber{#2}}\thmnote{\,--\,#3}}
\theoremstyle{mytheoremstyle}
\newmdtheoremenv[linewidth=1pt,backgroundcolor=shallowGreen,linecolor=deepGreen,leftmargin=0pt,innerleftmargin=20pt,innerrightmargin=20pt,]{theorem}{Theorem}[section]
\theoremstyle{mytheoremstyle}
\newmdtheoremenv[linewidth=1pt,backgroundcolor=shallowBlue,linecolor=deepBlue,leftmargin=0pt,innerleftmargin=20pt,innerrightmargin=20pt,]{definition}{Definition}[section]
\theoremstyle{myproblemstyle}
\newmdtheoremenv[linecolor=black,leftmargin=0pt,innerleftmargin=10pt,innerrightmargin=10pt,]{problem}{Problem}[section]
%%%%%%%%%%%%%%%%%%%%%%%%%%%%%%%%%%%%%%%%%%%%%%%%%%%%%%%%%%%%%%%%%%%%%%%%%%%%%%%

%%%%%%%%%%%%%%%%%%%%%%%%%%%%%%% Plotting Settings %%%%%%%%%%%%%%%%%%%%%%%%%%%%%
\usepgfplotslibrary{colorbrewer}
\pgfplotsset{width=8cm,compat=1.9}
%%%%%%%%%%%%%%%%%%%%%%%%%%%%%%%%%%%%%%%%%%%%%%%%%%%%%%%%%%%%%%%%%%%%%%%%%%%%%%%

%%%%%%%%%%%%%%%%%%%%%%%%%%%%%%% Title & Author %%%%%%%%%%%%%%%%%%%%%%%%%%%%%%%%
\title{Informe sobre FPGA y CPLD}
\author{Cortesini Perez Luciano Tomas}
\date{}
%%%%%%%%%%%%%%%%%%%%%%%%%%%%%%%%%%%%%%%%%%%%%%%%%%%%%%%%%%%%%%%%%%%%%%%%%%%%%%%

\begin{document}
\maketitle

\section{Introducción}
Este informe analiza la diferencia entre FPGA y CPLD, define las siglas principales, menciona los
fabricantes destacados y explica cómo se interpretan estos dispositivos a partir del material del
\href{https://www.youtube.com/watch?v=trBiZkW6POk}{video propuesto}.

\section{Definición y características de FPGA}
Un \textbf{FPGA} es un dispositivo lógico programable que contiene bloques lógicos configurables (CLB),
matrices de interconexión y puertos de E/S (GPIO).  
Los FPGAs permiten implementar funciones digitales complejas y pueden ser reprogramados después de
su fabricación. Generalmente ofrecen una gran cantidad de puertas lógicas, memoria interna y módulos
especializados para aritmética y procesamiento de señales.

\section{Siglas y significados}
En el video se identifican las siguientes siglas:
\begin{itemize}
    \item \textbf{PAL}: Programmable Array Logic, una familia de dispositivos programables simples.
    \item \textbf{PLA}: Programmable Logic Array, una matriz lógica más flexible que puede implementarse como combinaciones de puertas.
    \item \textbf{CPLD}: Complex Programmable Logic Device.
    \item \textbf{FPGA}: Field Programmable Gate Array.
    \item \textbf{SSI}: Stacked Silicon Interconnect, tecnología de interconexión entre bloques de silicio.
    \item \textbf{CLB}: Configurable Logic Block, el bloque básico de lógica en un FPGA.
    \item \textbf{PLL}: Phase Locked Loop, circuito que genera señales de reloj estables y sincronizadas.
    \item \textbf{MMCM}: Mixed Mode Clock Manager, unidad de manejo de relojes usada en FPGAs para generar y ajustar señales de reloj.
    \item \textbf{GPIO}: General Purpose Input Output, pines de propósito general para entrada y salida digital.
\end{itemize}

\section{Fabricantes principales}
Los fabricantes mencionados en el video son:
\begin{itemize}
    \item \textbf{Xilinx}.
    \item \textbf{Altera}.
\end{itemize}
Xilinx y Altera son empresas clave en el mercado de la lógica programable. Xilinx es conocida por
sus familias de FPGA y CPLD, y Altera, ahora parte de Intel, es reconocida por sus FPGAs y soluciones
de diseño integrado.

\section{Comparativa entre FPGA y CPLD}
Un \textbf{FPGA} y un \textbf{CPLD} son dos tipos de dispositivos lógicos programables que se usan
en distintos tipos de diseño, aunque comparten algunos principios.

\subsection*{Diferencias principales}
\begin{itemize}
    \item Los FPGAs son más adecuados para diseños grandes y complejos, con gran cantidad de puertas lógicas, bloques de memoria y recursos especializados.
    \item Los CPLDs suelen ser más compactos y están pensados para funciones de control, arranque y lógica de menor complejidad.
    \item La mayoría de los CPLDs incluyen memoria de configuración no volátil en chip, por lo que pueden iniciar inmediatamente al energizar el sistema.
    \item Los FPGAs, en cambio, suelen requerir cargar su configuración desde una memoria externa al arrancar, aunque tienen mayor capacidad interna.
    \item La distinción principal entre las arquitecturas de dispositivos FPGA y CPLD es que los CPLD están basados internamente en una colección de PLD acompañados por una estructura de interconexión programable, mientras que los FPGA usan bloques lógicos.
\end{itemize}

\subsection*{Aspectos compartidos}
\begin{itemize}
    \item Ambos dispositivos permiten implementar lógica digital programable y se basan en elementos lógicos configurables.
\end{itemize}

\subsection*{Uso típico}
\begin{itemize}
    \item Un CPLD suele usarse para gestión de arranque, inicialización de buses o tareas que requieren disponibilidad inmediata al encender.
    \item Un FPGA suele usarse para procesamiento digital complejo, aplicaciones de alto rendimiento y diseños que necesitan reconfiguración.
\end{itemize}

La diferencia más evidente entre un CPLD grande y un FPGA pequeño es la memoria no volátil del CPLD.  
Esta característica permite utilizar al CPLD como un dispositivo de "boot loader", cargando la
configuración de otros chips (como un FPGA) o inicializando señales críticas antes de que otros
componentes estén listos.

\section{Interpretación de estos dispositivos}
Los dispositivos lógicos programables son chips electrónicos que permiten crear circuitos
digitales de manera flexible. En lugar de tener una función fija grabada en el hardware, estos
dispositivos se pueden programar para realizar diferentes tareas lógicas, como procesar señales,
controlar motores o manejar datos.

Estos dispositivos sirven para simplificar el desarrollo de sistemas electrónicos. Antes, diseñar
un circuito personalizado requería fabricar un chip nuevo, lo que era costoso y lento. Con los
dispositivos lógicos programables, los ingenieros pueden probar ideas rápidamente, cambiar la
lógica sin alterar el hardware físico y reducir el tiempo de desarrollo. Se usan en una amplia
variedad de aplicaciones, desde controles industriales hasta dispositivos de consumo, permitiendo
innovaciones en tecnología sin grandes inversiones iniciales.

\section{Fuentes}
\begin{itemize}
    \item \href{https://en.wikipedia.org/wiki/Complex_programmable_logic_device}{Wikipedia: Complex Programmable Logic Device}
    \item \href{https://en.wikipedia.org/wiki/Field-programmable_gate_array}{Wikipedia: Field Programmable Gate Array}
    \item \href{https://www.youtube.com/watch?v=trBiZkW6POk}{Video YouTube}
\end{itemize}

\end{document}
